\cleardoublepage
\chapter{INSTRUMEN EVALUASI AHLI}
\label{lamp:instrumen}

Lampiran ini memuat instrumen yang digunakan pada evaluasi ahli sebagaimana hasilnya dilaporkan pada Subbab~\ref{sec:bab6-evaluasi-integrasi-cdss-ahli}. Butir
pertanyaan disalin apa adanya dari formulir yang diisi responden sehingga penilaian pada bab evaluasi dapat diperiksa terhadap pertanyaan yang benar-benar diajukan.
Identitas responden tidak disertakan pada lampiran ini maupun pada berkas statistik yang dibangkitkan.

Dua skala digunakan pada instrumen ini. Butir skenario, rekomendasi, dan aspek integrasi memakai skala 1 sampai 4, sedangkan butir \textit{System Usability Scale}
memakai skala 1 sampai 5 sesuai instrumen aslinya \parencite{brooke1996}. Perbedaan tersebut dipertahankan agar skor SUS dapat dihitung dengan algoritma bakunya.

\section{Pernyataan Persetujuan}
\label{lamp:instrumen-consent}

Formulir dibuka dengan pernyataan berikut, dan responden hanya dapat melanjutkan setelah menyatakan kesediaannya.

\begin{quote}
Dengan melanjutkan pengisian formulir ini, saya menyatakan bahwa saya telah membaca dan memahami informasi mengenai evaluasi sistem serta bersedia memberikan
tanggapan untuk keperluan penelitian Tugas Akhir.
\end{quote}

\section{Profil Responden}
\label{lamp:instrumen-profil}

Bagian profil dikumpulkan untuk menempatkan penilaian pada konteks pengalaman responden, bukan untuk menyaring peserta. Seluruh butir bersifat wajib kecuali
bidang okupasi dan keterangan pengalaman CDSS. Butir-butirnya ditunjukkan pada Tabel~\ref{tab:lampC-profil}.

\begingroup
\setlength{\LTcapwidth}{\textwidth}
\begin{longtable}{@{}>{\raggedright\arraybackslash}p{6.400cm}>{\raggedright\arraybackslash}p{7.178cm}@{}}
\caption{Butir profil responden pada instrumen evaluasi ahli.}
\label{tab:lampC-profil} \\
\toprule
\textbf{Pertanyaan} & \textbf{Pilihan jawaban} \\
\midrule
\endfirsthead

\multicolumn{2}{@{}l}{\small\emph{Tabel \thetable{} (lanjutan)}} \\
\toprule
\textbf{Pertanyaan} & \textbf{Pilihan jawaban} \\
\midrule
\endhead

\bottomrule
\endlastfoot

Apa profesi Anda? & Dokter Spesialis; Dokter; Dokter Muda; Perawat; Tenaga kesehatan lainnya; Mahasiswa kedokteran/kesehatan; Lainnya \\

Bidang okupasi & Isian bebas \\

Apakah Anda pernah menggunakan sistem CDSS atau aplikasi kesehatan digital sebelumnya? & Sudah; Belum \\

Jika pernah, sebutkan & Isian bebas \\

Apakah Anda pernah menangani pasien diabetes yang menggunakan CGM? & Ya; Jarang; Belum Pernah \\
\end{longtable}
\endgroup

Distribusi jawaban pada butir-butir tersebut dilaporkan pada Tabel~\ref{tab:tab6.16_karakteristik_responden}. Butir bidang okupasi tidak diisi seluruh responden,
sehingga cacahnya dilaporkan terpisah dari cacah responden keseluruhan.

\section{Penilaian Skenario Penggunaan}
\label{lamp:instrumen-skenario}

Butir pada bagian ini dinilai pada skala 1 sampai 4, dengan 1 menyatakan sangat tidak setuju dan 4 menyatakan sangat setuju.

Tiap skenario merupakan tugas terpandu yang dijalankan responden pada purwarupa, bukan naskah keluaran yang disajikan siap baca. Responden diminta memilih pasien
tertentu, menetapkan sumber pengukuran, menjalankan penilaian klinis, lalu menilai keluaran yang muncul. Rincian keenam skenario disajikan pada
Tabel~\ref{tab:lampC-skenario-tugas}.

\begingroup
\setlength{\LTcapwidth}{\textwidth}
\small
\captionsetup{font=normalsize}
\setlength{\tabcolsep}{4pt}
\begin{longtable}{@{}>{\raggedright\arraybackslash}p{1.300cm}>{\raggedright\arraybackslash}p{3.300cm}>{\raggedright\arraybackslash}p{1.500cm}>{\raggedright\arraybackslash}p{2.100cm}>{\raggedright\arraybackslash}p{4.675cm}@{}}
\caption{Skenario tugas yang dijalankan responden pada evaluasi ahli.}
\label{tab:lampC-skenario-tugas} \\
\toprule
\textbf{Kode} & \textbf{Nama skenario} & \textbf{Pasien} & \textbf{Sumber glukosa} & \textbf{Yang diamati responden} \\
\midrule
\endfirsthead

\multicolumn{5}{@{}l}{\small\emph{Tabel \thetable{} (lanjutan)}} \\
\toprule
\textbf{Kode} & \textbf{Nama skenario} & \textbf{Pasien} & \textbf{Sumber glukosa} & \textbf{Yang diamati responden} \\
\midrule
\endhead

\bottomrule
\endlastfoot

1 & Pemantauan dengan CGM & P001 & CGM & Prediksi pada horizon 30 dan 60 menit beserta label kondisi dan daftar rujukan \\

2 & Pemantauan dengan finger-stick & P007 & Finger-stick & Prediksi pada horizon sekitar empat jam dan cara sistem menyampaikan keterbatasan data \\

3A & Advisory pada kondisi stabil & P004 & CGM & Hubungan kondisi, prediksi, dan advisory, serta nomor halaman pada tiap rujukan \\

3B & Advisory pada hipoglikemia & P003 & CGM & Kesesuaian label kondisi dengan arah glukosa, dan pembahasan ketidakpastian pada advisory \\

3C & Advisory pada hiperglikemia & P002 & CGM & Kesesuaian rujukan dengan kondisi terprediksi, bukan dengan kondisi saat ini \\

4 & Data tidak memadai & P008 & CGM & Penjelasan mengapa prediksi tidak tersedia, pencatatan logbook baru, dan pendaftaran pasien baru \\
\end{longtable}
\endgroup

Skenario 3A sampai 3C sengaja menempatkan tiga kondisi glukosa yang berbeda pada pasien yang berbeda agar penilaian tidak bertumpu pada satu keadaan saja. Skenario
3C secara khusus meminta responden memeriksa apakah rujukan yang ditampilkan sesuai dengan kondisi yang diprediksi dan bukan dengan kondisi yang sedang berlaku,
yaitu perilaku yang menjadi kebaruan penelitian ini. Skenario 4 menguji jalur ketika data tidak memenuhi syarat prediksi, sehingga yang dinilai adalah cara sistem
menolak memberikan prediksi, bukan mutu prediksinya. Butir penilaian untuk tiap skenario ditunjukkan pada Tabel~\ref{tab:lampC-skenario}.


\begingroup
\setlength{\LTcapwidth}{\textwidth}
\begin{longtable}{@{}>{\raggedright\arraybackslash}p{2.300cm}>{\raggedright\arraybackslash}p{11.278cm}@{}}
\caption{Butir penilaian skenario penggunaan pada skala 1--4.}
\label{tab:lampC-skenario} \\
\toprule
\textbf{Skenario} & \textbf{Pernyataan} \\
\midrule
\endfirsthead

\multicolumn{2}{@{}l}{\small\emph{Tabel \thetable{} (lanjutan)}} \\
\toprule
\textbf{Skenario} & \textbf{Pernyataan} \\
\midrule
\endhead

\bottomrule
\endlastfoot

Prediksi CGM & Prediksi dan status kondisi konsisten dengan data atau konteks pasien yang tersedia. \\

Prediksi finger-stick & Sistem tetap memberikan prediksi yang relevan sekaligus mengomunikasikan keterbatasan data finger-stick dengan jelas. \\

Penolakan prediksi & Sistem menjelaskan dengan jelas mengapa prediksi tidak tersedia dan membedakannya dari kondisi normal. \\

Komunikasi kepastian & Sistem menghindari kesan kepastian berlebihan meski data pasien terbatas. \\
\end{longtable}
\endgroup

Selain butir tersebut, tiga kondisi klinis dinilai melalui kisi penilaian yang sama. Kondisi 3A merupakan kondisi stabil, 3B hipoglikemia, dan 3C hiperglikemia.
Setiap kondisi dinilai pada empat kriteria, yaitu relevansi, kejelasan, kecukupan konteks, dan batasan keputusan, masing-masing pada skala 1 sampai 4.

\section{Penilaian Kesesuaian Rekomendasi}
\label{lamp:instrumen-rekomendasi}

Pada kondisi hipoglikemia dan hiperglikemia, responden memilih satu kategori berikut untuk menilai kesesuaian rekomendasi sistem terhadap pedoman klinis.

\begin{enumerate}
    \item Sesuai pedoman dan aman diikuti.
    \item Sesuai pedoman, tetapi perlu verifikasi dokter.
    \item Kurang sesuai pedoman, berpotensi berisiko.
    \item Tidak sesuai pedoman.
\end{enumerate}

\section{System Usability Scale}
\label{lamp:instrumen-sus}

Sepuluh butir pada Tabel~\ref{tab:lampC-sus} dinilai pada skala 1 sampai 5. Butir bernomor ganjil berarah positif dan butir bernomor genap berarah negatif,
sehingga rerata yang lebih tinggi pada butir genap justru menunjukkan kebergunaan yang lebih rendah.

\begingroup
\setlength{\LTcapwidth}{\textwidth}
\begin{longtable}{@{}>{\raggedleft\arraybackslash}p{1.200cm}>{\raggedright\arraybackslash}p{9.957cm}>{\raggedright\arraybackslash}p{2.000cm}@{}}
\caption{Sepuluh butir System Usability Scale pada skala 1--5.}
\label{tab:lampC-sus} \\
\toprule
\textbf{No.} & \textbf{Pernyataan} & \textbf{Arah} \\
\midrule
\endfirsthead

\multicolumn{3}{@{}l}{\small\emph{Tabel \thetable{} (lanjutan)}} \\
\toprule
\textbf{No.} & \textbf{Pernyataan} & \textbf{Arah} \\
\midrule
\endhead

\bottomrule
\endlastfoot

1 & Saya berpikir akan sering menggunakan sistem ini. & Positif \\

2 & Saya merasa sistem ini rumit tanpa alasan yang jelas. & Negatif \\

3 & Saya merasa sistem ini mudah digunakan. & Positif \\

4 & Saya merasa membutuhkan bantuan teknis untuk menggunakan sistem ini. & Negatif \\

5 & Saya merasa fungsi-fungsi dalam sistem ini terintegrasi dengan baik. & Positif \\

6 & Saya merasa terdapat terlalu banyak inkonsistensi dalam sistem ini. & Negatif \\

7 & Saya membayangkan kebanyakan orang akan cepat mempelajari sistem ini. & Positif \\

8 & Saya merasa sistem ini sangat merepotkan untuk digunakan. & Negatif \\

9 & Saya merasa percaya diri saat menggunakan sistem ini. & Positif \\

10 & Saya perlu mempelajari banyak hal sebelum bisa menggunakan sistem ini dengan baik. & Negatif \\
\end{longtable}
\endgroup

Skor SUS dihitung per responden mengikuti algoritma Brooke, yaitu nilai butir ganjil dikurangi satu, nilai butir genap dikurangkan dari lima, lalu seluruhnya
dijumlahkan dan dikalikan 2,5 sehingga menghasilkan skor 0 sampai 100. Rerata butir dan ringkasan skornya disajikan pada Tabel~\ref{tab:tab6.19_ringkasan_sus}.

\section{Aspek Integrasi CDSS}
\label{lamp:instrumen-integrasi}

Butir pada Tabel~\ref{tab:lampC-integrasi} dinilai pada skala 1 sampai 4.

\begingroup
\setlength{\LTcapwidth}{\textwidth}
\begin{longtable}{@{}>{\raggedright\arraybackslash}p{3.300cm}>{\raggedright\arraybackslash}p{10.278cm}@{}}
\caption{Butir aspek integrasi CDSS pada skala 1--4.}
\label{tab:lampC-integrasi} \\
\toprule
\textbf{Aspek} & \textbf{Pernyataan} \\
\midrule
\endfirsthead

\multicolumn{2}{@{}l}{\small\emph{Tabel \thetable{} (lanjutan)}} \\
\toprule
\textbf{Aspek} & \textbf{Pernyataan} \\
\midrule
\endhead

\bottomrule
\endlastfoot

Populasi sumber data & Karakteristik populasi sumber data, yaitu DMT1, pengguna pompa insulin, dua belas pasien, dan populasi Amerika Serikat, memadai untuk mendukung penerapan sistem pada pasien DMT1 Indonesia. \\

Otonomi dokter & Sistem menjaga otonomi dokter dalam pengambilan keputusan klinis dan tidak berpotensi bias terhadap kelompok pasien tertentu. \\

Consent dan keamanan data & Mekanisme persetujuan dan keamanan data pasien dalam alur sistem ini sudah memadai. \\

Dasar rekomendasi & Dasar atau alasan di balik rekomendasi sistem cukup jelas untuk membangun kepercayaan klinis. \\

Alur kerja & Alur kerja sistem ini realistis diterapkan dalam praktik klinis Anda sehari-hari di Indonesia. \\
\end{longtable}
\endgroup

Kelima aspek tersebut dilaporkan bersama butir skenario pada Tabel~\ref{tab:tab6.17_evaluasi_skenario_cdss}. Butir populasi sumber data sengaja menyebutkan
karakteristik dataset secara eksplisit agar responden menilai kesenjangan populasi berdasarkan informasi yang sama.

\section{Isian Bebas}
\label{lamp:instrumen-bebas}

Empat butir isian bebas disediakan, dan seluruhnya bersifat opsional kecuali butir terakhir. Butir tersebut adalah catatan untuk skenario 3A sampai 3C, catatan
untuk skenario penolakan prediksi dan komunikasi kepastian, gap populasi yang dinilai paling signifikan, serta kritik, saran, atau masukan tambahan yang mencakup
hambatan penerapan, informasi yang kurang, dan risiko yang perlu dipertimbangkan. Rangkuman tematik jawaban pada butir-butir tersebut disajikan pada
Subbab~\ref{subsec:bab6-masukan-terbuka-ahli}.
